\section{Type Inference}\label{sec:type-inference-new}

In this section, we describe our type inference algorithm, where the goal is to annotate each subexpression of a program with its most precise type. Specifically, we describe our type inference algorithm for expressions with float type $\float[B;V]$, showing how \Slice{} populates the cut set $B$ and value set $V$ for each float-typed expression, as these refinements will later determine how and whether the expression can be discretized. Type inference for all other types is standard and follows the classic Hindley-Milner type inference algorithm. For our type inference algorithm on float-typed expressions, we follow a standard two-phase approach: the first phase traverses the program once to generate a set of constraints over cut and value sets; the second phase solves these constraints to produce a fully annotated program.

\subsection{Constraint Generation}

We use a constraint language as shown in Figure~\ref{fig:constraints-grammar} to express requirements over type structure and float refinements. We adopt the framework presented in~\cite{Odersky1999ConstrainedTypes}, introducing a new judgment of form $\Gamma \vdash e : \tau \mid C$, which reads as \emph{in the environment $\Gamma$, $e$ has type $\tau$ and produces the set of constraints $C$}. In this judgment, $\Gamma$ and $e$ are inputs to the type inference algorithm, while $\tau$ and $C$ are outputs—the algorithm infers the type $\tau$ for expression $e$ and generates the constraints $C$ that must be satisfied for this typing to be valid. The inference rules formalize what constraint $C$ must be generated (and later solved) for a given expression $e$ to have inferred type $\tau$. Figure~\ref{fig:type-inference} gives the constraint generation rules for \Slice{}. We focus on a small but representative subset of constructs in \Slice{}, but enough to capture the core features needed to demonstrate how contraints are generated over the full language.

The type inference rules closely mirror the typing rules in Figure~\ref{fig:typing-rules}, with two key differences. First, the inference rules generate a constraint set instead of having those constraints be part of the premises of the rules. Second, subtyping is handled in a syntax-directed way, where instead of a general subtyping rule, we bake in the subtyping relation into each rule.

\begin{figure}[!t]
    \begin{align*}
    C ::= &\; \tau_1 <: \tau_2  \text{\quad (subtyping)} \\
          &\mid \distinguishes{B}{V}  \text{\quad (distinguishability)} \\
          &\mid \answerslt{B}{V_1}{V_2}  \text{\quad (less-than answerability)} \\
          &\mid \answersleq{B}{V_1}{V_2}  \text{\quad (less-equal answerability)} \\
          &\mid \has{V}{c}  \text{\quad (value containment)} \\
          &\mid B_1 \sqsubseteq B_2  \text{\quad (cut set subset)} \\
          &\mid B_1 = B_2  \text{\quad (cut set equality)} \\
          &\mid V_1 \sqsubseteq V_2  \text{\quad (value set subset)} \\
          &\mid V_1 = V_2  \text{\quad (value set equality)} \\
          &\mid C_1, C_2   \text{\quad (conjunction)}
    \end{align*}
\caption{Syntax for our constraint language. The constraints directly correspond to the semantic requirements in our type system: subtyping constraints ensure type compatibility, $\distinguishes{B}{V}$ ensures that distribution parameters can be recovered after discretization, $\answerslt{B}{V_1}{V_2}$ and $\answersleq{B}{V_1}{V_2}$ ensure that comparison results can be determined after discretization, and $\has{V}{c}$ ensures that constants are properly tracked in value sets.}
\caption{Syntax for our constraint language. As we want \Slice{} to propogate information about comparisons and possible value ranges into float types, our syntax is over the cut sets $B \in \cutsets$ and value sets $V \in \valsets$. For the $\subtype{\tau_1}{\tau_2}$ constraint, we assume $\tau_1$ and $\tau_2$ are float types $\float[B_1;V_1]$ and $\float[B_2;V_2]$, respectively, and hence follow the \textsc{Sub-Float} relation in Figure~\ref{fig:subtyping}.}
\label{fig:constraints-grammar}
\end{figure}

\begin{figure}
    \begin{mathpar}
        \inferrule[\textsc{TFloat-Const}]
        {\ }
        {C \land \has{V}{c}, \Gamma \vdash c : \float[B;V]}

        \inferrule[\textsc{TVar}]
        {\ }
        {C, \Gamma, x: \tau \vdash x : \tau}
    
        \inferrule[\textsc{TContinuous-1}]
        {\ C, \Gamma \vdash e : \float[B_1; V_1]}
        {C \land \has{V}{\top} \land \distinguishes{B_1}{V_1}, \Gamma \vdash \exponential(e) : \float[B; V]}

        % \inferrule[\textsc{TContinuous-2}]
        % {\ C_1, \Gamma \vdash e_1 : \float[B_1; V_1] \\
        %    C_2, \Gamma \vdash e_2 : \float[B_2; V_2] }
        % {C_1 \land C_2 \land V = \top \land B_1 =\; \cleq{V_1} \land B_2 =\; \cleq{V_2}, \Gamma \vdash \uniform(e_1, e_2) : \float[B; \top]}
    
        \inferrule[\textsc{TIf}]
        {C_1, \Gamma \vdash e_1 : \bool \\
         C_2, \Gamma \vdash e_2 : \tau_2 \\
         C_3, \Gamma \vdash e_3 : \tau_3}
        {C_1 \land C_2 \land C_3 \land \subtype{\tau_1}{\tau_4} \land \subtype{\tau_2}{\tau_4}, \Gamma \vdash \ifkw \; e_1 \; \thenkw \; e_2 \; \elsekw \; e_3 : \tau_4}

        \inferrule[\textsc{TLess}]
        {C_1, \Gamma \vdash e_1 : \float[B; V_1] \and C_2, \Gamma \vdash e_2 : \float[B; V_2]}
        {C_1 \land C_2 \land \answerslt{B}{V_1}{V_2}, \Gamma \vdash e_1 < e_2 : \bool}

        \inferrule[\textsc{TLessEq}]
        {C_1, \Gamma \vdash e_1 : \float[B; V_1] \and C_2, \Gamma \vdash e_2 : \float[B; V_2]}
        {C_1 \land C_2 \land \answersleq{B}{V_1}{V_2}, \Gamma \vdash e_1 \leq e_2 : \bool}

        \inferrule[\textsc{Let}]
        {C_1, \Gamma \vdash e_1 : \tau_1 \\
         C_2, \Gamma, x: \tau_1 \vdash e_2 : \tau_2}
        {C_1 \land C_2, \Gamma \vdash \letkw \; x = e_1 \; \inkw \; e_2 : \tau_2}

    \end{mathpar}
\caption{Constraint generation rules for type inference for float-typed expressions.}
    \label{fig:type-inference}
\end{figure}

\begin{example}
    To further illustrate how our constraint generation rules work, we walk through an example below. Consider the program, omitting the irrelevant types such as \bool:

    \begin{lstlisting}[aboveskip=1em,belowskip=1em,escapechar=!,mathescape]
        let x = if flip()
                then exponential(0.5)
                else exponential(0.8)
                in
        let y = if x < 1.5
                then 1.8
                else 0.3
                in
        x <= y
    \end{lstlisting}
    
    Let us walk through the constraint generation for this program. We annotate each subexpression with type variables for the cut sets ($B_i$) and value sets ($V_i$):
    
    \begin{lstlisting}[aboveskip=1em,belowskip=1em,escapechar=!,mathescape]
        let x = if flip()
                then exponential(0.5 : !{\ftbn{$B_1$}{$V_1$}}!) : !{\ftbn{$B_2$}{$V_2$}}!
                else exponential(0.8 : !{\ftbn{$B_3$}{$V_3$}}!) : !{\ftbn{$B_4$}{$V_4$}}!
                : !{\ftbn{$B_5$}{$V_5$}}! in
        let y = if (x : !{\ftbn{$B_5$}{$V_5$}}!) < (1.5 : !{\ftbn{$B_6$}{$V_6$}}!)
                then (1.8 : !{\ftbn{$B_7$}{$V_7$}}!)
                else (0.3 : !{\ftbn{$B_8$}{$V_8$}}!)
                : !{\ftbn{$B_9$}{$V_9$}}! in
        (x : !{\ftbn{$B_5$}{$V_5$}}!) <= (y : !{\ftbn{$B_9$}{$V_9$}}!)
    \end{lstlisting}
    
    Applying our constraint generation rules, we obtain the following constraints:
    
    \begin{enumerate}
        \item From the constants using \textsc{TFloat-Const}:
        \begin{itemize}
            \item $\has{V_1}{0.5}$ (from literal 0.5)
            \item $\has{V_3}{0.8}$ (from literal 0.8)
            \item $\has{V_6}{1.5}$ (from literal 1.5)
            \item $\has{V_7}{1.8}$ (from literal 1.8)
            \item $\has{V_8}{0.3}$ (from literal 0.3)
        \end{itemize}
        
        \item From the exponential distributions using \textsc{TContinuous-1}:
        \begin{itemize}
            \item $\distinguishes{B_1}{V_1}, V_2 = \top$ (from \texttt{exponential(0.5)})
            \item $\distinguishes{B_3}{V_3}, V_4 = \top$ (from \texttt{exponential(0.8)})
        \end{itemize}
        
        \item From the first \texttt{if} expression using \textsc{TIf}:
        \begin{itemize}
            \item $\subtype{\float[B_2; V_2]}{\float[B_5; V_5]}$
            \item $\subtype{\float[B_4; V_4]}{\float[B_5; V_5]}$
        \end{itemize}
        
        \item From the comparison \texttt{x < 1.5} using \textsc{TLess}:
        \begin{itemize}
            \item $\answerslt{B_{10}}{V_5}{V_6}$
            \item $B_{10} = B_5, B_{10} = B_6$ (shared cut set requirement)
        \end{itemize}
        
        \item From the second \texttt{if} expression using \textsc{TIf}:
        \begin{itemize}
            \item $\subtype{\float[B_7; V_7]}{\float[B_9; V_9]}$
            \item $\subtype{\float[B_8; V_8]}{\float[B_9; V_9]}$
        \end{itemize}
        
        \item From the final comparison \texttt{x <= y} using \textsc{TLessEq}:
        \begin{itemize}
            \item $\answersleq{B_{11}}{V_5}{V_9}$
            \item $B_{11} = B_5, B_{11} = B_9$ (shared cut set requirement)
        \end{itemize}
    \end{enumerate}
    
    The complete constraint set generated is:
    \begin{align*}
    C = &\; \has{V_1}{0.5}, \has{V_3}{0.8}, \has{V_6}{1.5}, \has{V_7}{1.8}, \has{V_8}{0.3}, \\
        &\; V_2 = \top, V_4 = \top, \\
        &\; \distinguishes{B_1}{V_1}, \distinguishes{B_3}{V_3}, \\
        &\; \subtype{\float[B_2; V_2]}{\float[B_5; V_5]}, \subtype{\float[B_4; V_4]}{\float[B_5; V_5]}, \\
        &\; \subtype{\float[B_7; V_7]}{\float[B_9; V_9]}, \subtype{\float[B_8; V_8]}{\float[B_9; V_9]}, \\
        &\; \answerslt{B_{10}}{V_5}{V_6}, B_{10} = B_5, B_{10} = B_6, \\
        &\; \answersleq{B_{11}}{V_5}{V_9}, B_{11} = B_5, B_{11} = B_9
    \end{align*}
\end{example}

\subsection{Constraint Solving}
Once constraint generation completes, we obtain a system of constraints over cut sets $B$ and value sets $V$ for each float-typed expression in the program. The goal of the constraint solving phase is to assign concrete sets to each $B$ and $V$ that satisfy all generated constraints. Solving proceeds by interpreting the constraints as equations or inclusions over sets of floats, and applying a standard fixpoint algorithm to propagate equalities and bounds through the constraint graph.

The most common constraint forms involve equality ($C_1 = C_2$) and least upper bound operations ($V_1 \latticesup V_2 = V_3$), which correspond to set equality and union, respectively. For instance, a constraint of the form $V_1 \latticesup V_2 = V_3$ is solved by setting $V_3$ to the union of $V_1$ and $V_2$, while constraints such as $B = \clt{V}$ arising from comparisons are handled by computing the cut set of the corresponding value set. We solve constraints incrementally, propagating information until reaching a stable fixed point. For example, if we learn that $V_1 = \{0.3\}$ and $V_2 = \{1.8\}$, then the constraint $V_1 \latticesup V_2 = V_3$ allows us to conclude $V_3 = \{0.3, 1.8\}$. This, in turn, may determine $B_3 = \clt{V_3} = {\clt{0.3}, \clt{1.8}}$, which might then resolve other constraints downstream. The solving phase terminates when no further changes can be made to any $B$ or $V$ set.

\paragraph{Preserving observable output for partial discretization.} \katherine{todo}

\begin{example}
    From the previous example, we obtained the following constraints:

    \begin{enumerate}
        \item $B_1 = B_2 = B_3$ and $V_1 \latticesup V_2 = V_3$ 
        \item $B_4 = B_3$ and $B_4 = B_4\; \latticesup \clt{V_4}$
        \item $B_5 = B_6 = B_7$ and $V_5 \latticesup V_6 = V_7$ 
        \item $V_1 = V_2 = \top$, $V_4 = \{1.5\}$, $V_5 = \{1.8\}$, $V_6 = \{0.3\}$
        \item $B_3 = B_7$ and $B_3 = B_3\; \latticesup \cleq{V_7}$
    \end{enumerate}

\end{example}
